\documentclass[10pt,a4paper,twocolumn]{article}

\input{pisika.dat}

%  Editorial staff will uncomment the next line
% \input{staff.hed}


\begin{document}

%--------------------------------------------------------------------------
%  fill in the paper's title, author(s), and corresponding institutions
%--------------------------------------------------------------------------
\providecommand{\ShortAuthorList}[0]{A.~M.~Surname, B.~D.~Suffix Jr., C.~G. Suffix III} % use "A.~M.~Surname \textit{et al}." for more than three authors.
\title{More than one Author with different Affiliations}
\author[1]{Author M. Surname\thanks{Corresponding author: amsurname@university.edu}}
\author[2]{Bauthor D. Suffix Jr.}
\author[2]{Coauthor G. Suffix III}
\affil[1]{Department of Physics, DD University}
\affil[2]{Department of Science, XX University}

\date{\dateline{10 June 2016}}

\begin{abstract}
\noindent
%---------------------------------------------------------------------------
%               Include abstract and keywords here
%---------------------------------------------------------------------------
Perspiciatis unde omnis iste natus error sit voluptatem accusantium doloremque laudantium, totam rem aperiam, eaque ipsa quae ab illo inventore veritatis et quasi architecto beatae vitae dicta sunt explicabo. Nemo enim ipsam voluptatem quia voluptas sit aspernatur aut odit aut fugit, sed quia consequuntur magni dolores eos qui ratione voluptatem sequi nesciunt. Neque porro quisquam est, qui dolorem ipsum quia dolor sit amet, consectetur, adipisci velit, sed quia non numquam eius modi tempora incidunt ut labore et dolore magnam aliquam quaerat voluptatem. Ut enim ad minima veniam, quis nostrum exercitationem ullam corporis suscipit laboriosam, nisi ut aliquid ex ea commodi consequatur? Quis autem vel eum iure reprehenderit qui in ea voluptate velit esse quam nihil molestiae consequatur, vel illum qui dolorem eum fugiat quo voluptas nulla pariatur?

\keywords{keyword 1, keyword 2}

\DOI{} % do not delete this line
\end{abstract}

\maketitle
\thispagestyle{titlestyle}



%---------------------------------------------------------------------------
%               the main text of your paper begins here
%---------------------------------------------------------------------------
\section{Introduction}\label{sec:intro}
This is the template for articles submitted to the journal Pisika. Using the included BibTeX template one can cite articles \cite{articlekey}, arXiv preprints \cite{preprintkey}, books \cite{bookkey}, and conference proceedings \cite{proceedingskey}. It is preferred to insert citations before punctuation \cite{articlekey,preprintkey}. Sections, figures, and tables are referred to by the usual \textbackslash{\ttfamily{ref}}\{...\} command (as seen in Section \ref{sec:intro}).

Use the equation environment
\begin{verbatim}
\begin{equation}\label{eq:law2}
  \vec{F} = m \vec{a}.
\end{equation}	 
\end{verbatim}
to display equations
\begin{equation}\label{eq:law2}
		\vec{F} = m \vec{a} + \nu \bm\sigma.
\end{equation}
The {\ttfamily bm} package allows the use of bold face Greek letters: \textbackslash {\ttfamily bm}\textbackslash {\ttfamily sigma}. Equations are referred in text by the command \textbackslash{\ttfamily eqref}\{{\ttfamily eq:law2}\} as we do here \eqref{eq:law2}.
Longer equations may use the align environment
\begin{verbatim}
\begin{align}
\vec{F} & = \vec{F}_1 + \vec{F}_2 
          + \vec{F}_3 \nonumber\\
&\quad + \vec{F}_4 + \vec{F}_5 
          + \vec{F}_6 \nonumber\\
&\qquad + \vec{F}_7 + \vec{F}_8 + \vec{F}_9.
\end{align}
\end{verbatim}
to get
\begin{align}
\vec{F} & = \vec{F}_1 +  \vec{F}_2 
          +  \vec{F}_3  \nonumber \\
 &\quad + \vec{F}_4 +  \vec{F}_5 
          +  \vec{F}_6  \nonumber \\
 &\qquad + \vec{F}_7 + \vec{F}_8 + \vec{F}_9.
\end{align}

Ed ut perspiciatis unde omnis iste natus error sit voluptatem accusantium doloremque laudantium, totam rem aperiam, eaque ipsa quae ab illo inventore veritatis et quasi architecto beatae vitae dicta sunt explicabo. Nemo enim ipsam voluptatem quia voluptas sit aspernatur aut odit aut fugit, sed quia consequuntur magni dolores eos qui ratione voluptatem sequi nesciunt. Neque porro quisquam est, qui dolorem ipsum quia dolor sit amet, consectetur, adipisci velit, sed quia non numquam eius modi tempora incidunt ut labore et dolore magnam aliquam quaerat voluptatem. Ut enim ad minima veniam, quis nostrum exercitationem ullam corporis suscipit laboriosam, nisi ut aliquid ex ea commodi consequatur? Quis autem vel eum iure reprehenderit qui in ea voluptate velit esse quam nihil molestiae consequatur, vel illum qui dolorem eum fugiat quo voluptas nulla pariatur? ed ut perspiciatis unde omnis iste natus error sit voluptatem accusantium doloremque laudantium, totam rem aperiam, eaque ipsa quae ab illo inventore veritatis et quasi architecto beatae vitae dicta sunt explicabo. Nemo enim ipsam voluptatem quia voluptas sit aspernatur aut odit aut fugit, sed quia consequuntur magni dolores eos qui ratione voluptatem sequi nesciunt. Neque porro quisquam est, qui dolorem ipsum quia dolor sit amet, consectetur, adipisci velit, sed quia non numquam eius modi tempora incidunt ut labore et dolore magnam aliquam quaerat voluptatem. Ut enim ad minima veniam, quis nostrum exercitationem ullam ed ut perspiciatis unde omnis iste natus error sit voluptatem accusantium doloremque laudantium, totam rem aperiam, eaque ipsa quae ab illo inventore veritatis et quasi architecto beatae vitae dicta sunt explicabo. Nemo enim ipsam voluptatem quia voluptas sit aspernatur aut odit aut fugit, sed quia consequuntur magni dolores eos qui ratione voluptatem sequi nesciunt. 

\section{Figures}
Figures are inserted into the document by the environment
\begin{verbatim}
\begin{figure}[tb]
\centering
\includegraphics[width = ...]{figure1.pdf}
\caption{...}\label{fig:plot}
\end{figure}	
\end{verbatim}
shows how to insert Figure~\ref{fig:plot}. The option [{\ttfamily tb}] tells the compiler to place the figure at the top (first priority) or bottom (second priority) of the page. The figure can be scaled using the [{\ttfamily width}=...] option. Usually the width is specified as a fraction,  (say 0.6) of \textbackslash{\ttfamily linewidth}. This template can use *.pdf, *.png, *.jpg figures.

\begin{figure}[tb]
\centering
\includegraphics[width=0.6\linewidth]{figure1.pdf}
\caption{This is an example of a single column-wide figure. Figure captions are preferably self-contained and in vector format (versus raster format).}
\label{fig:plot}
\end{figure}

\begin{figure*}[t]
\centering
\includegraphics[width=0.8\linewidth]{widefigure.png}
\caption{This is an example of a figure spanning two columns. It will appear only on the top of pages. Please use these wide figures sparingly.}
\label{fig:wideplot}
\end{figure*}

Ed ut perspiciatis unde omnis iste natus error sit voluptatem accusantium doloremque laudantium, totam rem aperiam, eaque ipsa quae ab illo inventore veritatis et quasi architecto beatae vitae dicta sunt explicabo. Nemo enim ipsam voluptatem quia voluptas sit aspernatur aut odit aut fugit, sed quia consequuntur magni dolores eos qui ratione voluptatem sequi nesciunt. Neque porro quisquam est, qui dolorem ipsum quia dolor sit amet, consectetur, adipisci velit, sed quia non numquam eius modi tempora incidunt ut labore et dolore magnam aliquam quaerat voluptatem. Ut enim ad minima veniam, quis nostrum exercitationem ullam corporis suscipit laboriosam, nisi ut aliquid ex ea commodi consequatur? Quis autem vel eum iure reprehenderit qui in ea voluptate velit esse quam nihil molestiae consequatur, vel illum qui dolorem eum fugiat quo voluptas nulla pariatur?

Ed ut perspiciatis unde omnis iste natus error sit voluptatem accusantium doloremque laudantium, totam rem aperiam, eaque ipsa quae ab illo inventore veritatis et quasi architecto beatae vitae dicta sunt explicabo. Nemo enim ipsam voluptatem quia voluptas sit aspernatur aut odit aut fugit, sed quia consequuntur magni dolores eos qui ratione voluptatem sequi nesciunt. Neque porro quisquam est, qui dolorem ipsum quia dolor sit amet, consectetur, adipisci velit, sed quia non numquam eius modi tempora incidunt ut labore et dolore magnam aliquam quaerat voluptatem. Ut enim ad minima veniam, quis nostrum exercitationem ullam corporis suscipit laboriosam, nisi ut aliquid ex ea commodi consequatur? Quis autem vel eum iure reprehenderit qui in ea voluptate velit esse quam nihil molestiae consequatur?

\begin{table}[t]
\centering % center-align tables within a column
\caption{This is an example of a single column table. Captions are preferably self-contained and placed above the table. Columns may be left-, center-, decimal marker-, or right-aligned.}\label{table:label1}
\begin{tabular}{l c d{2.3} r}
\toprule
Example & Count 	1			& \multicolumn{1}{c}{Count 2} & Total\\
\midrule
A 						& 2.345 					& 5.435 																		& 16.78(3)\\
B 						& 3.0 							& 12.0																				& 43.2(5)\\
\bottomrule
\end{tabular}
\end{table}


\section{Tables}
In the sample Table~\ref{table:label1} the environment
\begin{verbatim} 
\begin{table}
\centering 
\caption{...}
 \begin{tabular}{l c d{2.3} r}
   \toprule
   ... & ... & ... & ... \\
   \midrule
   ... & ... & ... & ... \\
   \bottomrule
 \end{tabular}
\end{table}
\end{verbatim} 
constructs a four-column table with the first column left-aligned, second column centered, third column aligned on the decimail point with 2 integer digits and 3 decimal places, and the fourth column right-aligned. Ampersands {\ttfamily \&} separate columns while double backslashes \textbf{\textbackslash\textbackslash} \ start a new row.

Ed ut perspiciatis unde omnis iste natus error sit voluptatem accusantium doloremque laudantium, totam rem aperiam, eaque ipsa quae ab illo inventore veritatis et quasi architecto beatae vitae dicta sunt explicabo. Nemo enim ipsam voluptatem quia voluptas sit aspernatur aut odit aut fugit, sed quia consequuntur magni dolores eos qui ratione voluptatem sequi nesciunt. Neque porro quisquam est, qui dolorem ipsum quia dolor sit amet, consectetur, adipisci velit, sed quia non numquam eius modi tempora incidunt ut labore et dolore magnam aliquam quaerat voluptatem. Ut enim ad minima veniam, quis nostrum exercitationem ullam corporis suscipit laboriosam, nisi ut aliquid ex ea commodi consequatur? Quis autem vel eum iure reprehenderit qui in ea voluptate velit esse quam nihil molestiae consequatur, vel illum qui dolorem eum fugiat quo voluptas nulla pariatur?


Ed ut perspiciatis unde omnis iste natus error sit voluptatem accusantium doloremque laudantium, totam rem aperiam, eaque ipsa quae ab illo inventore veritatis et quasi architecto beatae vitae dicta sunt explicabo. Nemo enim ipsam voluptatem quia voluptas sit aspernatur aut odit aut fugit, sed quia consequuntur magni dolores eos qui ratione voluptatem sequi nesciunt. Neque porro quisquam est, qui dolorem ipsum quia dolor sit amet, consectetur, adipisci velit, sed quia non numquam eius modi tempora incidunt ut labore et dolore magnam aliquam quaerat voluptatem. Ut enim ad minima veniam, quis nostrum exercitationem ullam corporis suscipit laboriosam, nisi ut aliquid ex ea commodi consequatur? 

\section{Using BibTeX for the References}
The included BibTeX file {\ttfamily bibfile}.{\ttfamily bib} contains several sample entries. BibTeX formatted citations are easily downloadable from search engines. Please use standard journal abbreviations. As much as possible, do include the {\ttfamily DOI}=\{{\ttfamily 10.XXXX}/{\ttfamily XXXX}\}{\ttfamily ,} entry so that references are automatically hyperlinked (\textbf{do not include} the leading {\ttfamily http}:\textbackslash\textbackslash{\ttfamily dx.doi.org}\textbackslash).
\begin{table*}
  \centering
	\caption{This is an example of a table spanning two columns. It uses the \textbackslash{\ttfamily begin}\{{\ttfamily table*}\} ... \textbackslash {\ttfamily end}\{{\ttfamily table*}\} environment and is always placed on the top of pages. Please use these wide tables sparingly.} \label{table:wide}
  \begin{tabular}{lcccccccc}
	  \toprule
    1 & 2 & 3 & 2 & 3 & stretchstretchstretchstretch  stretchstretchstretchstretch & 3 & 2 & 3 \\       
		4 & 5 & 6 & 2 & 3 & stretchstretchstretchstretch  stretchstretchstretchstretch & 3 & 2 & 3 \\        
		7 & 8 & 9 & 2 & 3 & stretchstretchstretchstretch  stretchstretchstretchstretch & 3 & 2 & 3 \\
		\bottomrule
  \end{tabular}
\end{table*}

For a journal article, the recommended fields to be filled in are
\begin{verbatim}
@article{articlekey,
title={...},
author={Full Surname Sr., First N. and ...},
journal={J. Abbrev.},
volume={33},
pages={555},
year={2016},
doi={10.1000/doi.num.1000}
}
\end{verbatim}
Each author is separated by the word {\ttfamily and}. Suffixes are attached to the author last name.

Numquam eius modi tempora incidunt ut labore et dolore magnam aliquam quaerat voluptatem. Ut enim ad minima veniam, quis nostrum exercitationem ullam corporis suscipit laboriosam, nisi ut aliquid ex ea commodi consequatur? Quis autem vel eum iure reprehenderit qui in ea voluptate velit esse quam nihil molestiae consequatur, vel illum qui dolorem eum fugiat quo voluptas nulla pariatur?




ed ut perspiciatis unde omnis iste natus error sit voluptatem accusantium doloremque laudantium, totam rem aperiam, eaque ipsa quae ab illo inventore veritatis et quasi architecto beatae vitae dicta sunt explicabo. Nemo enim ipsam voluptatem quia voluptas sit aspernatur aut odit aut fugit, sed quia consequuntur magni dolores eos qui ratione voluptatem sequi nesciunt. Neque porro quisquam est, qui dolorem ipsum quia dolor sit amet, consectetur, adipisci velit, sed quia non numquam eius modi tempora incidunt ut labore et dolore magnam aliquam quaerat voluptatem. Ut enim ad minima veniam, quis nostrum exercitationem ullam corporis suscipit laboriosam, nisi ut aliquid ex ea commodi consequatur? Quis autem vel eum iure reprehenderit qui in ea voluptate velit esse quam nihil molestiae consequatur, vel illum qui dolorem eum fugiat quo voluptas nulla pariatur?


\section{Discussion}
Ed ut perspiciatis unde omnis iste natus error sit voluptatem accusantium doloremque laudantium, totam rem aperiam, eaque ipsa quae ab illo inventore veritatis et quasi architecto beatae vitae dicta sunt explicabo. Nemo enim ipsam voluptatem quia voluptas sit aspernatur aut odit aut fugit, sed quia consequuntur magni dolores eos qui ratione voluptatem sequi nesciunt. Neque porro quisquam est, qui dolorem ipsum quia dolor sit amet, consectetur, adipisci velit, sed quia non numquam eius modi tempora incidunt ut labore et dolore magnam aliquam quaerat voluptatem. Ut enim ad minima veniam, quis nostrum exercitationem ullam corporis suscipit laboriosam, nisi ut aliquid ex ea commodi consequatur? Quis autem vel eum iure reprehenderit qui in ea voluptate velit esse quam nihil molestiae consequatur, vel illum qui dolorem eum fugiat quo voluptas nulla pariatur?


Ed ut perspiciatis unde omnis iste natus error sit voluptatem accusantium doloremque laudantium, totam rem aperiam, eaque ipsa quae ab illo inventore veritatis et quasi architecto beatae vitae dicta sunt explicabo. Nemo enim ipsam voluptatem quia voluptas sit aspernatur aut odit aut fugit, sed quia consequuntur magni dolores eos qui ratione voluptatem sequi nesciunt. Neque porro quisquam est, qui dolorem ipsum quia dolor sit amet, consectetur, adipisci velit, sed quia non numquam eius modi tempora incidunt ut labore et dolore magnam aliquam quaerat voluptatem. Ut enim ad minima veniam, quis nostrum exercitationem ullam corporis suscipit laboriosam, nisi ut aliquid ex ea commodi consequatur? Quis autem vel eum iure reprehenderit qui in ea voluptate velit esse quam nihil molestiae consequatur?



Ed ut perspiciatis unde omnis iste natus error sit voluptatem accusantium doloremque laudantium, totam rem aperiam, eaque ipsa quae ab illo inventore veritatis et quasi architecto beatae vitae dicta sunt explicabo. Nemo enim ipsam voluptatem quia voluptas sit aspernatur aut odit aut fugit, sed quia consequuntur magni dolores eos qui ratione voluptatem sequi nesciunt. Neque porro quisquam est, qui dolorem ipsum quia dolor sit amet, consectetur, adipisci velit, sed quia non numquam eius modi tempora incidunt ut labore et dolore magnam aliquam quaerat voluptatem. Ut enim ad minima veniam, quis nostrum exercitationem ullam corporis suscipit laboriosam, nisi ut aliquid ex ea commodi consequatur? Quis autem vel eum iure reprehenderit qui in ea voluptate velit esse quam nihil molestiae consequatur, vel illum qui dolorem eum fugiat quo voluptas nulla pariatur?


Ed ut perspiciatis unde omnis iste natus error sit voluptatem accusantium doloremque laudantium, totam rem aperiam, eaque ipsa quae ab illo inventore veritatis et quasi architecto beatae vitae dicta sunt explicabo. Nemo enim ipsam voluptatem quia voluptas sit aspernatur aut odit aut fugit, sed quia consequuntur magni dolores eos qui ratione voluptatem sequi nesciunt. Neque porro quisquam est, qui dolorem ipsum quia dolor sit amet, consectetur, adipisci velit, sed quia non numquam eius modi tempora incidunt ut labore et dolore magnam aliquam quaerat voluptatem. Ut enim ad minima veniam, quis nostrum exercitationem ullam corporis suscipit laboriosam, nisi ut aliquid ex ea commodi consequatur? Quis autem vel eum iure reprehenderit qui in ea voluptate velit esse quam nihil molestiae consequatur, vel illum qui dolorem eum fugiat quo voluptas nulla pariatur?


\section{Conclusions}
Here is the Summary or Conclusions section.

\section*{Acknowledgments}
Here are the acknowledgments. Note the asterisk \textbackslash{\ttfamily section*}\{{\ttfamily Acknowledgments}\} that signifies that this section is unnumbered.

% Please use pisikabst.bst. You may your own *.bib file.
\bibliographystyle{pisikabst}
\bibliography{bibfile}



\end{document}